\section{Discussion}\label{sec:discussion}

\paragraph{Granger precedence is not causation.}
Granger causality tests whether past values of one series improve prediction of another beyond its own history; it does not establish true causation~\cite{granger1969investigating,shojaie2022granger}.
Bivariate Granger tests cannot exclude confounders: an unobserved process---such as a latent distributional shift in the training data---could independently drive both token entropy and diversity metrics, producing spurious lead-lag associations even when no direct relationship exists.
Our multi-method consensus framework inherits this limitation in full.
Throughout this paper, ``Granger precedence'' denotes predictive improvement in a time-series forecasting sense, not a causal mechanism.
This distinction is why we scope our central claim to ``token entropy Granger-precedes diversity collapse'' rather than asserting that token entropy \emph{causes} diversity collapse.

\paragraph{Calibration scope and generalization.}
All FPR and FNR calibration results reported in this work are specific to our experimental setup: Pythia-410M~\cite{biderman2023pythia}, a WikiText-103 subset, $T{=}10$ generation cycles, and five contamination levels ($\alpha \in \{0.0, 0.25, 0.50, 0.75, 1.0\}$).
Generalization to other model architectures, datasets, or training protocols requires independent validation.
The consensus framework's structure---multi-method agreement combined with multi-seed replication---is intended to be applicable beyond this specific setting, but the operating thresholds ($K{\geq}4/5$ methods, $S{\geq}3/4$ seeds) are empirically calibrated and may require re-tuning under different conditions.
For instance, architectures with different tokenizer granularity or training dynamics could shift the balance between false positive and false negative rates at these thresholds.
Our ongoing GPT-2-medium experiments represent a first step toward cross-architecture validation, though results are not yet available.

\paragraph{FPR/FNR calibration asymmetry.}
The FPR and FNR estimates in our calibration analysis differ in their coverage of the seed dimension.
FPR is validated on the method dimension only, using a single seed at the null-control condition ($\alpha{=}0.0$).
FNR is validated on the joint method${\times}$seed dimension, aggregating across four seeds at full contamination ($\alpha{=}1.0$).
This asymmetry means our FPR estimate of 0\% may be optimistic: if seed-level variance in false positive rates is non-negligible, the true FPR under the consensus rule could exceed zero.
Full joint calibration---requiring multi-seed null-control experiments at $\alpha{=}0.0$---remains an open direction for future work.

\paragraph{Short time-series constraint.}
With $T{=}10$ generation cycles (11 data points per series), our time-series analyses operate near the lower bound of what standard methods require.
This series length is typical for iterative training studies, where each generation involves a complete training run, but it severely limits statistical power.
The Toda--Yamamoto procedure's elevated FNR (33.3\%, missing 2 of 6 directional pairs) is attributable to insufficient residual degrees of freedom at this series length, not to a fundamental deficiency in the method itself.
This finding aligns with \citet{boettiger2012limits}, who demonstrated that early-warning signals have limited detection power in short time series.
Longer generation chains would improve power for all methods but impose substantial computational cost; this trade-off between statistical power and compute budget is inherent to iterative training research and cannot be resolved by methodological improvements alone.
However, observation window sensitivity analysis provides reassurance: the $4/5$ consensus classification stabilizes at $T{\geq}5$, well below our observation window of $T{=}10$, with identical results from $T{=}5$ through $T{=}10$.
The Toda--Yamamoto procedure produces valid test statistics at all window lengths ($T \in \{5, \ldots, 10\}$), but with unstable sensitivity---3/6 to 6/6 pairs reach significance depending on window length---confirming that its elevated FNR reflects limited residual degrees of freedom rather than absence of signal.

\paragraph{Diversity--calibration dissociation.}
Diversity metrics and calibration metrics respond to synthetic contamination in qualitatively opposite directions across the $\alpha$ spectrum.
Diversity collapse (distinct-$n$ decline) is most severe at low contamination ($\alpha{=}0.25$: distinct-3 declines by 59.9\%) and attenuates at higher levels ($\alpha{=}1.0$: $-30.5\%$), while calibration degradation (ECE) increases monotonically with $\alpha$ ($+51.7\%$ at $\alpha{=}0.0$ to $+3782\%$ at $\alpha{=}1.0$).
This dissociation underscores the need for multi-method consensus: the regime in which a metric degrades most severely ($\alpha{=}0.25$ for diversity) does not coincide with the regime in which its temporal lead is most detectable ($\alpha{\geq}0.75$), and a single metric family cannot capture the full picture of iterative training degradation.

\paragraph{Limitations.}
Several limitations constrain the generalizability of our findings.
First, our primary experiments use Pythia-410M; preliminary cross-architecture validation with GPT-2-medium uses batch size~32 (versus~4 for Pythia) for GPU efficiency, exposing 8$\times$ more data per training step, which may accelerate collapse dynamics without altering the qualitative pattern (token entropy leading diversity decline is observed in both architectures).
Second, training data is drawn from a single domain (a WikiText-103 subset), and domain-specific distributional properties may influence the coupling between canary and downstream metrics.
Third, we lack an intervention experiment: we observe that token entropy Granger-precedes diversity collapse at $\alpha{=}1.0$ (with preliminary single-seed evidence at $\alpha{=}0.75$) but do not test whether acting on canary signals can prevent or delay collapse onset.
Fourth, our experiments use replace-only data accumulation; accumulate-mode training~\cite{dey2024breaking,kazdan2025collapse} may exhibit qualitatively different canary dynamics.
Finally, our canary claims are scoped to distinct-$n$ diversity metrics~\cite{li2016diversity}---not to model collapse broadly~\cite{schaeffer2025position}---as the perplexity channel proved spurious ($F{\approx}0$, $p{=}0.99$ after differencing).
This scoping is theoretically motivated: \citet{cao2025entropycal} prove a fundamental calibration-diversity tradeoff in language models---all current calibration methods reduce entropy at the cost of diversity---suggesting that calibration and diversity are coupled but not identical, which may explain why token entropy can Granger-precede diversity collapse without leading perplexity.

\paragraph{Future work.}
Several directions extend this work.
Multivariate VAR models would enable testing for confounders and multi-directional Granger relationships beyond our current bivariate framework.
Cross-architecture validation is needed: our GPT-2-medium experiments are in progress, and evaluation on larger-scale models would further test generalizability.
Online monitoring---integrating canary metrics into training pipelines for real-time collapse detection---is a natural applied extension.
Accumulate-mode experiments~\cite{dey2024breaking} would test whether canary dynamics differ from the replace-mode setting studied here.
Finally, connecting the empirically observed detectability boundary at $\alpha{\approx}0.50$--$0.75$ to theoretical phase transition analysis~\cite{truong2025entropy} could yield a principled explanation for why canary signals emerge at this contamination regime rather than at lower contamination levels.